% 与questions_ch文件配套；来源超链接见sources.tex。
\Answer{SUPG2-01}{同一OR与不同概率差}
\textbf{答案：D。} 设未暴露者概率为$p_0$，由$\{p_1/(1-p_1)\}/\{p_0/(1-p_0)\}=3$，得$p_1=3p_0/(1-p_0+3p_0)$。两层的暴露者概率分别为0.25和0.75，故概率差分别为0.15和0.25。共同OR约束的是优势尺度上的倍数，不能据此要求各层有相同概率差，也不能把基线概率直接乘3。
\par\nobreak\smallskip\nobreak{\footnotesize\color{Muted}\noindent 来源与改编：知识点配套。题干与设问为新编；参考依据\hyperref[src:G2-NOTE]{[7]}。\par}

\Answer{SUPG2-02}{病例对照抽样后的概率校正}
\textbf{答案：B。} 样本优势为$0.60/0.40=1.5$。结局依赖抽样使样本优势等于人群优势乘$s_1/s_0$，因此人群优势为$1.5/12=0.125$，概率为$0.125/(1+0.125)=1/9\approx0.1111$。等价地，从样本线性预测子中减去$\log12$后取逆logit。直接把样本概率除以12得到0.05并不正确；校正发生在优势或logit尺度。
\par\nobreak\smallskip\nobreak{\footnotesize\color{Muted}\noindent 来源与改编：知识点配套。题干与设问为新编；参考依据\hyperref[src:G2-NOTE]{[7]}。\par}

\Answer{SUPG2-03}{分组二项响应的事件方向}
\textbf{答案：E。} 交换成功与失败两列后，模型化的概率变为$1-p$。由于$\operatorname{logit}(1-p)=-\operatorname{logit}(p)$，同一设计矩阵下全部系数（含截距）均变号，拟合概率互为补数。分组二项似然及其饱和模型也相应交换成功、失败，最大似然值和残差偏差不变。此结论依赖两次拟合使用相同资料和模型空间；不能只把斜率或只把截距变号。
\par\nobreak\smallskip\nobreak{\footnotesize\color{Muted}\noindent 来源与改编：知识点配套。题干与设问为新编；参考依据\hyperref[src:G2-NOTE]{[7]}、\hyperref[src:G1-GLM]{[2]}。\par}

\Answer{SUPG2-04}{三种检验分别需要哪些拟合}
\textbf{答案：A。} 有效比分检验在$H_0$约束下拟合干扰参数，再用约束点的得分和扣除干扰参数影响后的信息构造统计量，因此只需约束模型。Wald检验使用完整模型的估计及协方差；似然比检验须分别极大化约束和非约束模型。比分检验不是把其余系数也全部置零；多变量Wald检验同样可以调整模型中的其他协变量。
\par\nobreak\smallskip\nobreak{\footnotesize\color{Muted}\noindent 来源与改编：知识点配套。题干与设问为新编；参考依据\hyperref[src:G2-NOTE]{[7]}、\hyperref[src:G2-ANOVA]{[10]}。\par}

\Answer{SUPG2-05}{1:2匹配组的条件概率}
\textbf{答案：C。} 给定组内病例数为1，$A$的条件概率为$\exp(\beta x_A)/\{\exp(\beta x_A)+\exp(\beta x_B)+\exp(\beta x_C)\}=2/(2+1+1)=0.5$。组截距在条件化时消去。$2/(1+2)$只适用于本例中有1名暴露者和1名未暴露者的两人比较，没有计入第二名对照；这里计算的也不是$A$在人群中的绝对发病概率。
\par\nobreak\smallskip\nobreak{\footnotesize\color{Muted}\noindent 来源与改编：知识点配套。题干与设问为新编；参考依据\hyperref[src:G2-NOTE]{[7]}、\hyperref[src:G2-CLOGIT]{[11]}。\par}

\Answer{SUPG2-06}{条件Logistic可估计的内容}
\textbf{答案：B。} I正确：同组公共项在条件概率的分子、分母中相消。II错误：若匹配因素$Z$在组内恒定，交互项的组内差仍为$Z_i(x_{Ai}-x_{Bi})$；当不同$Z$层均有足够组内暴露变异、设计矩阵可识别时，可以估计暴露效应是否随$Z$改变。III正确：条件概率以组内病例数为条件，组截距已消去，不能直接恢复人群绝对风险。
\par\nobreak\smallskip\nobreak{\footnotesize\color{Muted}\noindent 来源与改编：知识点配套。题干与设问为新编；参考依据\hyperref[src:G2-NOTE]{[7]}、\hyperref[src:G2-CLOGIT]{[11]}。\par}

\Answer{SUPG2-07}{二次项与指定区间的OR}
固定其他协变量，两点的线性预测子之差为\[\Delta\eta=\beta_1\Delta+\beta_2\{(x_0+\Delta)^2-x_0^2\}=\beta_1\Delta+\beta_2(2x_0\Delta+\Delta^2).\]因此$\mathrm{OR}=\exp(\Delta\eta)$。当$\Delta=1$时，OR为$\exp\{\beta_1+\beta_2(2x_0+1)\}$，随起点$x_0$变化；只解释$e^{\beta_1}$会漏掉二次项贡献。教学系数下，$0\to1$对应$\exp(0.40-0.05)=1.4191$，而$4\to5$对应$\exp(0.40-0.45)=0.9512$。这反映logit曲线在不同区间的变化方向可能不同，不宜用单个“每单位效应”概括全部取值。\par 若要检验整个$x$项，应检验$H_0:\beta_1=\beta_2=0$，在通常正则条件下用2个自由度的联合检验。只检验$\beta_1=0$并不能排除二次关联；若检验是否需要二次项，则另用1个自由度检验$\beta_2=0$。
\par\nobreak\smallskip\nobreak{\footnotesize\color{Muted}\noindent 来源与改编：知识点配套。题干与设问为新编；参考依据\hyperref[src:G2-NOTE]{[7]}。\par}

\Answer{SUPG2-08}{食管癌公开资料中的饮酒评分约束}
$M_0$、$M_s$、$M_f$分别有$1+5+3=9$、10和12个参数。饮酒整体检验为$H_0:b_1=b_2=b_3=0$，\[G^2=2[-98.695896-(-162.662323)]=127.9329,\qquad df=3,\]提示调整年龄和吸烟后饮酒组与病例状态有关。\par 评分模型要求$(b_1,b_2,b_3)=(\theta,2\theta,3\theta)$，相当于$b_2-2b_1=0$、$b_3-3b_1=0$两个独立约束。\[G_s^2=2[-98.695896-(-101.646905)]=5.9020,\qquad df=2.\]在0.05水平未拒绝这些约束，$P=0.0523$，但这一接近阈值的结果并不能证明线性设定成立。\par 两个非参照饮酒组的比较必须使用系数差及其协方差：\[\begin{aligned}\hat d&=\hat b_2-\hat b_1=0.546089,\\\mathrm{SE}(\hat d)&=\sqrt{0.08108937+0.06253113-2(0.04266069)}=0.241452,\\\widehat{\mathrm{OR}}&=e^{\hat d}=1.7265,\qquad95\%\,CI=\exp\{\hat d\pm1.96(0.241452)\}=(1.0756,2.7714).\end{aligned}\]这是同年龄、同吸烟组条件下的病例优势比。评分的1单位表示跨一个人为排列的类别；最高组还是开放区间。没有个体实际饮酒克数，不能把评分约束解释为每增加固定克数的剂量线性关系，更不能仅凭未拒绝原假设加以证明。
\par\nobreak\smallskip\nobreak{\footnotesize\color{Muted}\noindent 来源与改编：案例改编。材料及改编范围见\hyperref[src:G2-NOTE]{[7]}、\hyperref[src:G2-ESOPH]{[8]}、\hyperref[src:G2-ESOPH-DATA]{[9]}、\hyperref[src:G2-ANOVA]{[10]}。\par}

\Answer{SUPG2-09}{1:2匹配六种模式的完整条件似然}
一组中共$c+k$名暴露者，病例的条件概率为$r^c/\{(c+k)r+3-c-k\}$。合并各模式，\[L_c(\beta)=3^{-7}\frac{r^{19}}{(r+2)^{14}(2r+1)^{15}},\qquad r=e^\beta.\]因此$U(\beta)=19-14r/(r+2)-30r/(2r+1)$。令得分为0，整理得$20r^2-21r-38=0$；唯一正根为$\hat r=2$，故$\hat\beta=\log2=0.693147$。所有人都未暴露的4组和都暴露的3组各贡献$1/3$，其贡献与$\beta$无关；剩余$36-7=29$组提供暴露效应信息。\par 信息给出$\mathrm{SE}(\hat\beta)=1/\sqrt{5.9}=0.411693$，\[95\%\,CI_{\rm OR}=\exp\{0.693147\pm1.96(0.411693)\}=(0.8925,4.4820).\]将$\hat r=2$代回完整条件似然，$\ell_c(\hat\beta)=-38.070179$，故\[G^2=2[-38.070179-(-39.550042)]=2.9597,\qquad df=1,\quad P=0.0854.\]在0.05水平未拒绝无关联假设，但区间较宽，不能据此认定暴露毫无作用。以上Wald和LR推断都是渐近近似。OR=2描述匹配条件下的优势比较；组截距已消去，且设计只固定组内病例数，因此不能从该OR或条件分配概率直接恢复人群绝对发病概率。
\par\nobreak\smallskip\nobreak{\footnotesize\color{Muted}\noindent 来源与改编：知识点配套。题干与设问为新编；参考依据\hyperref[src:G2-NOTE]{[7]}、\hyperref[src:G2-CLOGIT]{[11]}。\par}

\Answer{SUPG2-10}{平均预测正确是否足够}
全样本观察到$15+5+25+35=80$个事件，模型期望事件数为$50(0.10+0.30+0.50+0.70)=80$，两者总体比例均为0.40，故残差总和确为0。但这一等式只约束整体平均水平：A组观察比例0.30而预测0.10，低估10个事件；B组观察比例0.10而预测0.30，高估10个事件，两项误差恰好抵消。C、D组的观察比例分别为0.50、0.70，与预测相同。\par 在含截距、未惩罚且有限MLE存在的普通二项logit模型中，训练样本的截距得分方程本来就可导致预测与观察事件总数相等；它从不保证每个风险层均校准良好。本题是独立验证资料，这个总量等式即使仍成立，也不能升级为逐组校准的证明。\par 应报告各组观察比例和相应不确定性，绘制观察比例对预测概率的校准图，检查A、B组是否存在系统偏离及其协变量构成，必要时重新评估函数形式或校准。组内样本比例本身也有抽样波动，不能只凭差值就断言总体偏差的精确大小。若调整模型，仍需另外验证；平均比例相等不能替代这些检查，也不能评价全部预测性能。
\par\nobreak\smallskip\nobreak{\footnotesize\color{Muted}\noindent 来源与改编：知识点配套。题干与设问为新编；参考依据\hyperref[src:G2-NOTE]{[7]}、\hyperref[src:G1-GLM]{[2]}。\par}
